\subsection{Operating Points and Lead-Time Evaluation}

The reference selector deploys the smallest finite threshold whose raw risk upper bound meets the requested budget. This is a deterministic operating-point convention rather than a guarantee of optimal detection. \Cref{sec:utility-selector} adds an optional validation-utility stage for settings in which the clinical cost of a false alarm makes the smallest feasible threshold a poor operating point. Neither a lower threshold nor a higher budget is assumed to improve every operational metric of a stateful policy.

\subsection{Utility-Ranked Risk Calibration}
\label{sec:utility-selector}

An alarm budget constrains volume but does not encode the relative cost of a missed event and a false alarm. Let $U_i(\lambda)$ be patient $i$'s contribution to the normalised 2019 utility at threshold $\lambda$, and let $\mathcal V$ be a validation partition disjoint from both calibration and test. For a calibration procedure that returns a feasible set $\mathcal F_B$ at budget $B$, define
\begin{equation}
\label{eq:utility-select}
\widehat\lambda_U=\arg\max_{\lambda\in\mathcal F_B}
\left[\bar U_{\mathcal V}(\lambda)-z_{1-\delta}
\frac{s_{\mathcal V}(\lambda)}{\sqrt{|\mathcal V|}}\right],
\end{equation}
where $\bar U_{\mathcal V}$ and $s_{\mathcal V}$ are the validation mean and standard deviation of the patient-level contributions, and $z_{1-\delta}$ is the one-sided normal quantile at level $\delta=0.10$. Ties are broken toward the lower validation alarm rate; if $\mathcal F_B$ is empty the selector returns the exact no-alarm fallback. The lower bound discourages a threshold whose apparent utility is driven by a few high-variance patients. It is not a finite-sample utility guarantee.

The risk interpretation depends on how $\mathcal F_B$ was constructed. With the blocked empirical-Bernstein certificate, every candidate in $\mathcal F_B$ is simultaneously certified on the calibration sample, so validation ranking cannot invalidate the high-probability risk statement. With CRC, the expected-risk guarantee is attached to the CRC selection rule itself; the utility-ranked threshold is a deployment variant, not a second certificate. We therefore report the CRC row and the utility-ranked CRC (U-CRC) row separately, and include the validation alarm rate so that a reader can see whether the utility stage consumed more of the budget than intended. This distinction is deliberate: a data-dependent choice among CRC candidates should not be presented as inheriting the CRC guarantee.

For an event observed during follow-up, define the detection window $\mathcal W_i=[\max(0,t_i^\star-W),t_i^\star)$. An event is detected if at least one alarm falls in that window, and its lead time is
\begin{equation}
\label{eq:lead}
\ell_i(\lambda)=t_i^\star-\min\{t\in\mathcal W_i:a_{it}(\lambda)=1\}.
\end{equation}
Lead time is undefined for undetected events and is represented as missing rather than zero. The reported median therefore describes detected events only and must always be accompanied by sensitivity and the number of eligible events. Events at admission, outside observed follow-up, or without an identifiable onset cannot supply an observed pre-onset lead time. A retrospective top-$k$ selection over a complete trajectory is noncausal and may serve only as a labelled oracle comparator, not as a deployable early-warning policy.

Budget sweeps compare methods under the same requested upper-risk constraint, not necessarily at identical realised alarm rates. The plotted frontier uses realised alarm rate and conditional median lead time, retaining only nondominated pairs. At equal alarm rates only the largest lead time is retained; at larger rates a point remains only when its lead time strictly improves. Because conditional lead time can improve by missing difficult cases, sensitivity must be reported alongside the frontier before any clinical interpretation.

\begin{algorithm}[tb]
\caption{Patient-blocked upper-confidence calibration of an alarm policy}
\label{alg:calibration}
\begin{algorithmic}[1]
\Require independent calibration patients; fixed scorer, grid $\Lambda$, structural parameters, valid $L_{\max}$; budget $B$; level $\alpha$
\For{$\lambda\in\Lambda$}
  \State replay the causal policy on each patient's observed follow-up
  \State compute $L_i(\lambda)$ and assert $L_i(\lambda)\in[0,L_{\max}]$
  \State compute raw $\widehat R_n(\lambda)$ and Hoeffding or empirical-Bernstein radius
\EndFor
\State $\mathcal F\leftarrow\{\lambda\in\Lambda:\widehat R_n(\lambda)+\epsilon(\lambda)\le B\}$
\State $\widehat\lambda\leftarrow\min\mathcal F$ if nonempty; otherwise $+\infty$
\State \Return threshold, raw certificate, and explicit fallback status
\end{algorithmic}
\end{algorithm}
